\documentclass[11pt]{article}
\usepackage[a4paper,margin=25mm]{geometry}
\usepackage{amsmath,amssymb,amsthm,mathtools}
\usepackage[T1]{fontenc}
\usepackage{lmodern,microtype}
\usepackage[colorlinks=true,linkcolor=blue,urlcolor=blue]{hyperref}
\newcommand{\Ric}{\operatorname{Ric}}
\newcommand{\tr}{\operatorname{tr}}
\newcommand{\Vol}{\operatorname{Vol}}
\newcommand{\dd}{\,d}
\newcommand{\EH}{\mathrm{EH}}
\newcommand{\fint}{\operatorname{av}}
\newtheorem{proposition}{Proposition}
\title{Ricci-defect variation at a rational-ball end\\\large Exact inner minimization and a cutoff-dependent comparison}
\author{Research continuation}
\date{11 October 2026}
\begin{document}
\maketitle
\noindent
The target remains a smooth closed global minimizer of
\[
 P(g)=\int_M|E_g|^2\dd V_g,\qquad E_g=\Ric_g-\tfrac14s_g g,
\]
which is not Bach-flat. This note computes two precise parts of the remaining
rational-ball smoothing problem. It does not establish smooth attainment.
The first calculation is an exact minimum for a prescribed quadratic
asymptotic jet. The second includes an entire cutoff annulus and shows why
an isolated negative finite part is not the insertion coefficient of a fixed
critical orbifold.

\section{The zero-energy core and its quadratic response}
Let
\[
 G=1-r^{-4},\qquad
 g_0=G^{-1}dr^2+r^2(\sigma_1^2+\sigma_2^2+G\sigma_3^2),\qquad r\ge1.
\]
The forms are normalized so that $\sum\sigma_i^2$ is the unit round
three-sphere metric. Take the usual free involution of $\EH$ covering the
antipodal map of the bolt. Its quotient is the ordinary oriented rational
ball $B_2$, with link $S^3/\mathbb Z_4$. Write
\[
 L=\Vol(S^3/\mathbb Z_4)=\pi^2/2.
\]
The marking and quotient action are fixed throughout. The metric is
Ricci-flat and anti-self-dual. More generally, an ALE scalar-flat K\"ahler
rational homology ball is Ricci-flat, and has zero mass
\cite{Wright,HeinLeBrun}. Thus the variable scalar-flat K\"ahler mass available
on other fillings is absent in this class.

Consider the biaxial variation
\[
 g_t=H^{-1}dr^2+r^2(\sigma_1^2+\sigma_2^2+F\sigma_3^2),\quad
 F=G+tu+O(t^2),\quad H=G+tv+O(t^2).
\]
At a fixed bolt, the linear smoothness conditions are
\[
 u(1)=v(1)=0,\qquad u_r'(1)+v_r'(1)=0.
\]
Higher order boundary corrections impose $F'(1)H'(1)=16$. They do not
change the quadratic $P$ coefficient on a fixed compact region, since
$E_{g_0}=0$. Put $x=r^2$ and set
\begin{align}
 u&=\frac{\nu-\mu}{6}x+\frac{\nu+\mu}{2x}
               -\frac{2\nu+\mu}{3x^2},\label{response}\\
 v&=u-\nu(x-x^{-1}).\nonumber
\end{align}
This is a smooth quotient-invariant linear response. In the curvature
operator convention $R^+=W^++sI/12$, its first variation satisfies
\[
 DR^+[h]=\operatorname{diag}(\nu,\nu,\mu),\qquad
 Ds[h]=4(2\nu+\mu),\qquad h=\dot g_0.
\]
The homogeneous quadratic jet is Einstein at the cone vertex:
$\Ric(p)=(2\nu+\mu)g(p)$. This restriction is stronger than mere
$\mathbb Z_4$ invariance; arbitrary trace-free Ricci jets are not included.

\begin{proposition}[Exact inner minimum]
For the quadratic jet specified by \eqref{response},
\[
 DE_{g_0}[h]=\mu o,\qquad
 o=r^{-4}\operatorname{diag}(1,-1,-1,1)
\]
in the $g_0$-orthonormal frame. In particular,
\begin{equation}
 \int_{B_2}|DE[h]|^2=\frac{\pi^2}{2}\mu^2.\label{square}
\end{equation}
If $k$ is a smooth quotient-invariant correction with
$k=O_2(r^{2-\delta})$, $\delta>0$, and $DE[k]\in L^2$, then
\begin{equation}
 \int_{B_2}|DE[h+k]|^2
 =\frac{\pi^2}{2}\mu^2+\int_{B_2}|DE[k]|^2.\label{minimum}
\end{equation}
Thus \eqref{square} is the minimum in this specified class, including
non-biaxial corrections.
\end{proposition}
\begin{proof}
Direct calculation of sectional curvatures gives
\[
 D\Ric[h]=(2\nu+\mu)I+
 \mu r^{-4}\operatorname{diag}(1,-1,-1,1).
\]
For example, with $x$ derivatives,
\begin{align*}
 K_{01}=K_{02}&=-H_x,&
 K_{12}&=(4-3F-H)/x,\\
 K_{13}=K_{23}&=(F-H)/x-HF_x/F,\\
 K_{03}&=-H(3F_x/F+2xF_{xx}/F-xF_x^2/F^2)
              -H_x(1+xF_x/F).
\end{align*}
Substitution of \eqref{response} proves the asserted Ricci variation.
Since $dV_{g_0}=r^3dr\,d\omega$,
\[
 \|o\|_2^2=4L\int_1^\infty r^{-5}dr=L.
\]
The tensor $o$ is the smooth transverse-traceless scale Jacobi field;
$(\nabla^*\nabla-2\mathring R)o=0$, as in the classical
Eguchi--Hanson deformation calculation \cite{Biquard}.
Pairing $DE[k]$ with $o$ removes the gauge and trace terms. Green's
identity gives zero: there is no bolt boundary, and the boundary at
radius $R$ is $O(R^{-\delta})$. Expanding the square proves
\eqref{minimum}.
\end{proof}

There is no logarithmic divergence in \eqref{square}. Under an inner
expansion $g_a=a^2(g_0+a^2h+O(a^4))$, it contributes
$a^4L\mu^2$ on expanding rescaled cores, subject to matching estimates.
It is not a complete closed-manifold expansion. When $\mu=0$, this
quadratic obstruction vanishes; it supplies no strict exclusion of that
degeneration.

\section{An angular mode: inner finite part versus the full cutoff}
Use
\[
 \sigma_1^2+\sigma_2^2=\tfrac14(d\theta^2+\sin^2\theta\,d\phi^2),
 \qquad \sigma_3=\tfrac12(d\psi+\cos\theta\,d\phi).
\]
A harmonic Gibbons--Hawking moment coordinate is
\[
 f=\sqrt{r^4-1}\,Y,\qquad Y=\sin\theta\cos\psi,
 \qquad f_0=r^2Y.
\]
Choose the free involution preserving this moment coordinate; $f$ then
descends to $B_2$. The angular average is $\fint Y^2=1/3$.
The vertical-angle term is essential: the radial equation is
\[
 r^{-3}(r^3GA')'-r^{-2}\{8+4(G^{-1}-1)\}A=0,
 \qquad A=\sqrt{r^4-1}.
\]
One obtains
\[
 \fint|df|^2=4r^2-\frac4{3r^2},\qquad
 \int_{r\le R}|\nabla^2 f|^2
   =L\left(4R^4-\frac83-\frac4{3R^4}\right)
\]
by the Ricci-flat Bochner identity. For $h_f=2fg_0$,
$DE[h_f]=-2\nabla^2f$, and hence
\begin{equation}
 \lim_{R\to\infty}\left\{
 \int_{r\le R}|DE[h_f]|^2-
 \int_{r\le R}|DE_{\mathrm{flat}}[2f_0g_{\mathrm{flat}}]|^2
 \right\}=-\frac{32L}{3}=-\frac{16\pi^2}{3}.\label{negative}
\end{equation}
At the same time $Ds[h_f]=0$ and $DW^+[h_f]=0$.
The finite part \eqref{negative} therefore cannot be identified with
the chiral bulk integral while discarding its boundary term.

This can be checked with an entire smooth cutoff, not merely the inner
ball. Let $\chi(\rho)=1$ near $\rho=0$ and $\chi=0$ for large $\rho$,
and define $f_R=\chi(r/R)f$ and $f_{0,R}=\chi(r/R)f_0$.
For compactly supported conformal changes of either Ricci-flat model,
Chern--Weil and the scalar transformation law give exactly
\[
 P(e^{2tu}g_0)=3\int(t\Delta u+t^2|du|^2)^2\dd V_{g_0}.
\]
Here $P(g_0)=0$, and the conformal change is identical outside a compact
set, so the boundary correction does not change.

\begin{proposition}[Complete cutoff comparison]
Let $[t^2]$ denote the coefficient of $t^2$, rather than the second
derivative. Then
\begin{align}
 \lim_{R\to\infty}[t^2]\{P(e^{2tf_R}g_0)
  -P(e^{2tf_{0,R}}g_{\mathrm{flat}})\}
  =-3L\int_0^\infty
      \bigl(\rho^3(\chi'')^2+3\rho(\chi')^2\bigr)\dd\rho.
 \label{cutoff}
\end{align}
The right side depends on the cutoff profile.
\end{proposition}
\begin{proof}
Write $B_0=\rho^2\chi''+7\rho\chi'$. The exact radial product rule gives
\[
 \Delta(\chi(r/R)f)
 =\left(AG\chi_{rr}+A(7/r+r^{-5})\chi_r\right)Y.
\]
On the fixed support of the derivatives of $\chi$, this equals
\[
 \left[B_0+R^{-4}\left(-\tfrac32\rho^{-2}\chi''
                         -\tfrac52\rho^{-3}\chi'\right)
       +O(R^{-8})\right]Y.
\]
Both models have radial volume $r^3dr\,d\omega$.
The limit of the difference of their squared-Laplacian integrals is
\[
 \frac L3\int\{-3\rho^3(\chi'')^2
             -26\rho^2\chi''\chi'-35\rho(\chi')^2\}\dd\rho
 =-L\int\{\rho^3(\chi'')^2+3\rho(\chi')^2\}\dd\rho.
\]
Multiplication by $3$ proves \eqref{cutoff}.
\end{proof}
With $z=\log\rho$ and $\psi(z)=\chi(e^z)$, the profile integral is
\[
 \int_{\mathbb R}\bigl((\psi'')^2+4(\psi')^2\bigr)\dd z.
\]
If $\psi(z)=\psi_1(z/T)$ is a transition of logarithmic width $T$,
this becomes $T^{-3}\|\psi_1''\|_2^2+4T^{-1}\|\psi_1'\|_2^2$.
It tends to zero as $T\to\infty$. In particular the isolated number in
\eqref{negative} is not the answer to the fixed-bulk insertion problem.
The backgrounds in \eqref{cutoff} are being varied with $\chi$ and $R$;
they have not been required to be closed $P$-critical metrics.

\section{The remaining closed-manifold comparison}
The inner minimization and the cutoff calculation answer distinct parts
of the same question. A closed $P$-critical orbifold has a fixed exterior
geometry. Its Euler--Lagrange equation removes the interior first-variation
pairing, but a singular ALE tail leaves a boundary concomitant. The
subleading inner tail, angular components, and higher exterior jets must
be matched before the finite $a^4$ coefficient has an invariant meaning.
Neither \eqref{square} nor \eqref{negative} has determined that coefficient.

There is a concrete reason not to set that boundary term to zero using
stationarity. The leading flat-coordinate Eguchi--Hanson tail is
\[
 h^\infty_{ij}=\frac{x_i x_j-(Jx)_i(Jx)_j}{r^6}
 =\tfrac18(\partial_{ij}-J_i{}^kJ_j{}^l\partial_{kl})r^{-2}.
\]
Its distributional linearized Ricci defect pairs with a trace-free,
divergence-free test tensor through the contraction
$J_i{}^kJ_j{}^l\partial_{kl}E_{ij}(0)$, whereas the flat principal
critical equation controls $\Delta E$. These are different contractions.
For example, take $Je_1=e_3$, $Je_2=e_4$ and
\[
 E=(x_1^2-x_2^2)\operatorname{diag}(0,0,1,-1).
\]
Then $\tr E=\operatorname{div}E=\Delta E=0$, but
$J_i{}^kJ_j{}^l\partial_{kl}E_{ij}=4$.
It is generated by the TT quartic metric jet
$h_4=-(x_1^4-x_2^4)\operatorname{diag}(0,0,1,-1)/6$ through
$DE[h_4]=-\Delta h_4/2=E$.
This is a statement about the principal jet equations, not a construction
of a global critical metric. It identifies information that those equations
alone do not eliminate.

Two existing tools do not supply this missing comparison. With tensor
Weyl norm, set $\mathcal B_t=\int|W|^2+t\int s^2$. Then
\[
 2P=\mathcal B_{1/6}-32\pi^2\chi(M).
\]
At a closed $\mathcal B_t$-critical metric with $t\ne0$,
\[
 \nabla P=(t-1/6)sE.
\]
Thus the non-Einstein gluing constructions at the different parameters
in \cite{GV} do not furnish $P$-critical metrics. Also, for an order-four
Ricci-flat ALE end, inversion gives
\[
 \bar g=\delta+O_2(\rho^4),\qquad W_{\bar g}(p)=0.
\]
The compactification-point Weyl interaction in \cite{MM} therefore
vanishes for this end. Higher asymptotic curvature interactions remain,
and any comparison for $P$ must include the scalar-square term as well.

The unresolved statement is a strict closed-manifold energy comparison
against the singular minimizing configurations that can actually occur,
together with control of the other compactness alternatives. No smooth
global non-Bach-flat minimizer has been produced in this continuation.

\begin{thebibliography}{9}
\bibitem{LeBrun} C. LeBrun, \emph{Curvature Functionals, Optimal Metrics,
and the Differential Topology of 4-Manifolds}, Section 11.
\href{https://arxiv.org/abs/math/0404251}{arXiv:math/0404251}.
\bibitem{Biquard} O. Biquard, \emph{D\'esingularisation de m\'etriques
d'Einstein. I}, Sections 1--3.
\href{https://arxiv.org/abs/1105.4904}{arXiv:1105.4904}.
\bibitem{Wright} E. Wright, \emph{Quotients of gravitational instantons},
Corollary C. \href{https://arxiv.org/abs/1102.2442}{arXiv:1102.2442}.
\bibitem{HeinLeBrun} H.-J. Hein and C. LeBrun, \emph{Mass in K\"ahler
geometry}. \href{https://arxiv.org/abs/1507.08885}{arXiv:1507.08885}.
\bibitem{GV} M. J. Gursky and J. A. Viaclovsky, \emph{Critical metrics
on connected sums of Einstein four-manifolds}.
\href{https://arxiv.org/abs/1303.0827}{arXiv:1303.0827}.
\bibitem{MM} A. Malchiodi and F. Malizia, \emph{Weyl energy and connected
sums of four-manifolds}, (1.6)--(1.7).
\href{https://arxiv.org/abs/2505.17752}{arXiv:2505.17752}.
\end{thebibliography}
\end{document}
